
\documentclass[../../ps_q.tex]{subfiles}

\begin{document}

\section{回路の諸々}

\subsection{金属電子論}
    \subfile{01_metal-electron-theory/ps_em_cx_01_q.tex}

    \vspace{2cm}
    \vfill

\subsection{非オーム抵抗}
    \subfile{02_non-ohmic-resistor/ps_em_cx_02_q.tex}

    \vspace{2cm}
    \vfill

\subsection{合成抵抗・回路の対称性}
    \subfile{03_equivalent-resistance/ps_em_cx_03_q.tex}

    \vspace{2cm}
    \vfill

\subsection{コンデンサの無限回つなぎかえ①}
    \subfile{04_infinite-switching-1/ps_em_cx_04_q.tex}

    \vspace{2cm}
    \vfill 

\subsection{コンデンサの無限回つなぎかえ②}
    \subfile{05_infinite-switching-2/ps_em_cx_05_q.tex}   

    \vspace{2cm}
    \vfill

\subsection{誘電体}
    \subfile{06_dielectric/ps_em_cx_06_q.tex}

    \vspace{2cm}
    \vfill

\subsection{※導線中の自由電子の運動}
    \subfile{07_add-free-electrons/ps_em_cx_07_q.tex}

    \vspace{2cm}
    \vfill

\subsection{※電球を含む回路}
    \subfile{08_add-bulb/ps_em_cx_08_q.tex}

    \vspace{2cm}
    \vfill

\subsection{※合成抵抗}
    \subfile{09_add-equivalent-resistance/ps_em_cx_09_q.tex}

    \vspace{2cm}
    \vfill

\subsection{※コンデンサのつなぎかえ}
    \subfile{10_add-switching/ps_em_cx_10_q.tex}

    \vspace{2cm}
    \vfill

\subsection{※誘電体を挿入したコンデンサ}
    \subfile{11_add-dielectric/ps_em_cx_11_q.tex}

    \vspace{2cm}
    \vfill

\end{document}
